\section{Power differentiation and a positive generator}
\label{sec:tangent}

We first differentiate the power parameter and then express the result as the action of a positive operator on logarithmic rates.
The gamma--Dirichlet representation supplies the derivative; the posterior identity and the bounded phase formula supply the operator.

\subsection{Differentiating the power parameter}

Throughout this subsection $U$ is a positive measure on $(0,\infty)$ with
\begin{equation}
 0<B=U((0,\infty))<\infty,\qquad
 \int\log(1+b^{-1})\,U(\dd b)<\infty.
 \label{eq:finite-admissible}
\end{equation}
Let $X$ have the zero-drift GGC law with Thorin measure $U$.
Thus
\[
 L(s)=\E e^{-sX}
 =\exp\left\{-\int\log(1+s/b)\,U(\dd b)\right\}.
\]
For $s>0$, let $X_s$ have the exponential tilt of $X$ defined in Lemma~\ref{lem:gamma-dirichlet}.
Recall the digamma function $\psi_0=\Gamma'/\Gamma$.
For a probability measure $P$ on $(0,\infty)$, put
\begin{equation}
 M_P(s)=\int\frac{P(\dd b)}{s+b},\qquad
 W_P(s)=M_P(s)+sM_P'(s)=\int\frac{b\,P(\dd b)}{(s+b)^2}.
 \label{eq:MW}
\end{equation}

\begin{lemma}[The power derivative]
\label{lem:tangent}
Let $U$, $B$, and $X$ be as above.
Let $P\sim\DP(U)$ and $G_B\sim\operatorname{Gamma}(B,1)$ be independent.
Then, for $s>0$,
\begin{equation}
 X_s\overset d=G_BM_P(s).
 \label{eq:tilt-gamma-dirichlet}
\end{equation}
For $s,q>0$, define
\[
 g_q(s)=-\partial_s\log\E e^{-sX^q},\qquad
 h_U(s)=\left.\partial_q g_q(s)\right|_{q=1},\qquad g(s)=g_1(s).
\]
The derivatives defining $g_q$ and $h_U$ exist, and expectations on the right-hand sides below are with respect to $P\sim\DP(U)$:
\begin{align}
 g(s)&=B\E M_P(s),\label{eq:g-dirichlet}\\
 h_U(s)&=B\E\left[
 W_P(s)\{\psi_0(B+1)+\log M_P(s)\}+sM_P'(s)
 \right],\label{eq:tangent}\\
 \frac{h_U(s)+g(s)}B
 &=\E\left[W_P(s)\{\psi_0(B+1)+1+\log M_P(s)\}\right].
 \label{eq:normalized-tangent}
\end{align}
Equivalently, $h_U(s)=\partial_s\{s\E[X_s\log X_s]\}$.
\end{lemma}

\begin{proof}
The tilted representation \eqref{eq:tilt-gamma-dirichlet} is \eqref{eq:gamma-dirichlet-tilt}.
We use it as an equality of distributions for each fixed $s$.

For $q$ in a compact subinterval of $(0,\infty)$, and $s$ in a compact positive interval, differentiating $e^{-s x^q}$ in $s$ and $q$ produces finite sums of bounded functions of $z=x^q$ of the form $z^r(\log z)^j e^{-sz}$ with $r>0$, multiplied by bounded powers of $q^{-1}$.
At $x=0$ these terms have the continuous value zero.
Differentiation under the original probability is therefore justified without a moment assumption on $X$.
Differentiating first in $q$ yields
\[
 \left.\partial_q\log\E e^{-sX^q}\right|_{q=1}
 =-s\E[X_s\log X_s],
\]
and hence $h_U(s)=\partial_s\{s\E[X_s\log X_s]\}$.

To evaluate the tilted moment, we differentiate the gamma integral with respect to its positive shape parameter and obtain $\E[G_B\log G_B]=B\psi_0(B+1)$.
Consequently
\[
 \E[X_s\log X_s]
 =B\E\left[M_P(s)\{\psi_0(B+1)+\log M_P(s)\}\right].
\]
It remains to differentiate this expression in $s$.
For every $P$ and integer $j\ge1$,
\[
 0<W_P(s)\leq M_P(s)\leq s^{-1},\qquad
 |M_P^{(j)}(s)|\leq j!s^{-j}M_P(s).
\]
On compact positive $s$-intervals, the needed derivatives of $M_P(s)\log M_P(s)$ are bounded by a constant times $M_P(s)(1+|\log M_P(s)|)$, uniformly bounded for $0<M_P(s)\leq s^{-1}$.
Thus we may differentiate this last expectation, which gives \eqref{eq:tangent}.
Equation~\eqref{eq:g-dirichlet} follows either from the Thorin representation or from the mean of the tilted law.
Adding it to \eqref{eq:tangent} and using $W_P=M_P+sM_P'$ proves \eqref{eq:normalized-tangent}.
\end{proof}

The lemma applies to any current zero-drift GGC law with a finite Thorin measure.
This is what will allow us to use the same identity at every time in the evolution.
